\chapter{La promesa \emph{plug-and-play}}
\label{cap:plug-and-play}

\begin{intuicion}
La regla \emph{plug-and-play} de GraphHunter dice: cada
optimización entregada (S/M/L) se activa por defecto al levantar
el binario. Sin flags a setear, sin config a tocar, sin toggles en
la UI. El usuario abre la app y todo lo que está \emph{shipped}
se nota inmediatamente. Esa restricción no es cosmética: hace que
el motor sea \emph{auditable}. Si después de un ship de
optimización el analista tuvo que setear
\texttt{GRAPHHUNTER\_X=1} para verla, la promesa se rompió ---la
optimización no está realmente entregada, está \emph{disponible}.
Este capítulo codifica la regla en cinco pasos concretos y discute
por qué es deliberada y no accidental.
\end{intuicion}

\section{La promesa, en una frase}

\begin{quote}
\itshape
``Ninguna optimización entregada requiere que el usuario cambie su
config.''
\end{quote}

\noindent
Variantes equivalentes:

\begin{itemize}
    \item Cada flag S/M/L que ships se setea \emph{default-on} al
    \emph{startup} del binary.
    \item Cada cache es transparente.
    \item Cada índice se construye solo (lazy build).
    \item Cada índice grande se spills a disco solo.
\end{itemize}

\noindent
El usuario abre la app, ve la pildora ``Perf: M+L1'' verde, y todo
funciona.

\section{Por qué la regla es deliberada}

Sin la regla, el motor cae en un patrón de fail-mode común en
software de seguridad: \emph{shipped pero no usado}. El equipo
agrega una optimización, la testea, la merge ---pero el usuario
nunca la activa porque:

\begin{itemize}
    \item No sabe que existe (no leyó el changelog).
    \item Sabe que existe pero no sabe cómo activarla.
    \item Sabe cómo activarla pero le da fiaca tocar config.
    \item Tocó config, anduvo bien, pero después actualizó el
    binary y el config volvió a defecto.
\end{itemize}

\noindent
Default-on resuelve los cuatro casos: el código nuevo se activa
sin acción del usuario.

\section{La regla en cinco pasos}

\subsection{Paso 1: default-on al startup}

Cada flag nuevo se setea en \filepath{apps/tauri/src-tauri/src/lib.rs}
y \filepath{apps/cli/src/main.rs} con:

\begin{lstlisting}[language=Rust]
for (key, default) in [
    ("GRAPHHUNTER_LFTJ_PLANNER", "1"),  // M1 structural K_4
    ("GRAPHHUNTER_K4_LFTJ", "1"),        // S3 name marker
    ("GRAPHHUNTER_YANNAKAKIS", "1"),     // L1 stage 2/3
] {
    if std::env::var(key).is_err() {
        unsafe { std::env::set_var(key, default) };
    }
}
\end{lstlisting}

\noindent
Override con \texttt{KEY=0} sigue funcionando para usuarios que
quieran el path legacy (debugging, A/B comparison).

\subsection{Paso 2: surface en UI}

Extender \texttt{cmd\_get\_perf\_status} para incluir el flag
nuevo, y la pildora del toolbar para reflejarlo. El analista debe
poder \emph{ver} qué optimizaciones están activas sin abrir
DevTools ni leer logs (capítulo~\ref{cap:frontend-instr}).

\subsection{Paso 3: default-off sólo con causa}

Hay flags que NO son default-on (durante un período). Justifica
sólo si:

\begin{itemize}
    \item El flag re-escribe el DFS hot loop (cambios estructurales
    grandes).
    \item Aún no hay parity tests aguantando \(\geq 1\) semana
    estable.
    \item Hay un riesgo conocido de regresión en algún workload.
\end{itemize}

\noindent
En esos casos, el flag queda \emph{opt-in} durante un \emph{soak
window}. Cuando los parity tests aguantan, se flipea a
default-on.

\subsection{Paso 4: doc obliga}

Cada flip default-off $\to$ default-on actualiza:

\begin{itemize}
    \item \filepath{08-env-vars.tex} (libro Especificación).
    \item \filepath{12-preguntas-frecuentes.tex} (libro Manual).
    \item \texttt{make all} para refrescar los 6 PDFs.
\end{itemize}

\noindent
Sin doc actualizado, el flip no se merge.

\subsection{Paso 5: capabilities de build}

Los \emph{Cargo features} (como \texttt{symmetry-nauty}) que
agregan funcionalidad opcional al binary se reflejan via
constantes \texttt{HAS\_NAUTY}-style en
\texttt{cmd\_get\_perf\_status}. El frontend lee la capability y
muestra el feature como ``disponible'' o ``desactivado en este
build'' ---sin que el analista tenga que abrir
\filepath{Cargo.toml}.

\section{Cuándo se rompe la promesa}

Casos reales en los que casi se rompió:

\begin{itemize}
    \item \textbf{LFTJ-K\(_4\) inicial (S3)}: requería
    \texttt{GRAPHHUNTER\_K4\_LFTJ=1} \emph{y} que el nombre de la
    hipótesis contuviera ``k4-clique''. Cumplía la cláusula 1
    (default-on) pero el name marker era una barrera. Se mantuvo
    por back-compat pero M1 (structural\_dispatch) lo dejó
    redundante.
    \item \textbf{Yannakakis L1}: el flag default-on. Pero hubo un
    período donde el matcher Yannakakis tenía un bug en hipótesis
    con predicados; el ``fix'' fue requerir
    \texttt{GRAPHHUNTER\_YANNAKAKIS=0} para esos casos. \emph{No}
    cumplía con la regla. Se revertió, se arregló el bug en el
    operador, y el flag se restauró a default-on.
\end{itemize}

\noindent
La regla obliga a \emph{arreglar el bug en el operador}, no a
\emph{ofrecer un opt-out al usuario}.

\section{Por qué \emph{auditable}}

Plug-and-play hace al motor auditable: el equipo de Base4 puede
demostrar a un auditor externo que el motor en producción usa las
optimizaciones publicadas. Sin la regla, demostrar eso requeriría
inspeccionar configs en cada deployment ---imposible a escala.

Con la regla: el binary mismo es la prueba. ``Si lo bajaste, lo
estás usando.''

\section{Resumen}

\begin{itemize}
    \item Plug-and-play: cada optimización entregada se activa
    por defecto al startup.
    \item Cinco pasos: default-on al startup, surface en UI,
    default-off sólo con causa, doc obliga, capabilities de build
    reflejadas.
    \item Justificada porque sin ella, optimizaciones quedan
    \emph{shipped pero no usadas}.
    \item Casos donde casi se rompió: LFTJ S3 (name marker),
    Yannakakis bug temporal. Se priorizó arreglar el código
    sobre ofrecer opt-out.
    \item Auditabilidad: el binary mismo es la prueba.
\end{itemize}

\begin{ejercicios}
\begin{enumerate}
    \item \starred{} Imaginá una nueva optimización: ``cache de
    NLF compartido entre hipótesis similares''. ¿Cumple los cinco
    pasos? Discutí qué necesita.

    \item ¿Por qué el override \texttt{KEY=0} sigue funcionando?
    ¿Es contradictorio con la promesa plug-and-play?

    \item \starred{} Diseñá un test automatizado que verifique
    ``ninguna optimización merge-eada requiere setear flag
    manualmente''. ¿Qué chequearía?

    \item Casos donde la regla \emph{no} aplica: optimizaciones
    que sí necesitan opt-in del usuario por su naturaleza
    (riesgosa, costosa). Pensá un ejemplo.

    \item Si una optimización tiene un \emph{soak window} de 2
    semanas, ¿qué métricas decidirían que ya está listo para
    flip a default-on?
\end{enumerate}
\end{ejercicios}
